\section{Method: \methodname{}}
\label{sec:method}

We present \methodname{} (Cross-Modal Compensated Drift LoRA), a parameter-efficient buffer-free framework for cross-domain multimodal CIL.
The method is designed around the structural prior identified in Section~\ref{sec:observation}: in independent-branch fusion architectures, the spectral branch is naturally stable across domains while spatial branches require adaptation.
Fig.~\ref{fig:overview} illustrates the overall architecture.

\input{figures/fig_framework}

\subsection{\backbonename{} Backbone}
\label{sec:backbone}

\backbonename{} (Spectral-Spatial Contrastive Mamba) is a tri-branch multimodal backbone that processes HSI+LiDAR input through three feature streams:

\begin{itemize}
    \item \textbf{Spectral branch}: A 1D Mamba selective state-space model~\cite{gu2024mamba} processes per-pixel spectral vectors of the HSI input, capturing material absorption signatures via efficient long-sequence modeling. Output: $\fspec \in \R^{d}$.
    \item \textbf{HSI-spatial branch}: After a modality-specific input projection, HSI patches are processed by a stack of 2D VSSBlocks (depthwise convolution + channel MLP with state-space mixing) that extract spatial texture patterns. Output: $\fhsi \in \R^{d}$.
    \item \textbf{LiDAR-spatial branch}: LiDAR elevation/intensity maps pass through an independent input projection and then the \emph{same} VSSBlock stack in a Siamese configuration (shared spatial parameters, independent inputs). Output: $\flid \in \R^{d}$.
\end{itemize}

A lightweight cross-fusion module exchanges information between the two spatial streams after the final VSSBlock via residual cross-modal projections, before global average pooling produces the branch-level features.
The three outputs are concatenated into $f = [\fspec;\, \fhsi;\, \flid] \in \R^{3d}$ for classification.
The architecture uses $d{=}64$, 4 VSSBlocks per spatial branch, and 2 Mamba layers in the spectral branch.\footnote{All reported results use the Mamba~\cite{gu2024mamba} and VMamba VSSBlock implementations. The released code includes fallback paths to standard Transformer blocks when these dependencies are unavailable; fallback paths were not used for any reported experiment.}

The key architectural property is \emph{spectral--spatial independence}: the spectral branch shares no parameters with the spatial processing pipeline.
Section~\ref{sec:observation} shows that this separation preserves the drift asymmetry that \methodname{} exploits---the spectral branch naturally maintains stable representations while spatial branches absorb domain-specific variation, creating the structural prior that enables asymmetric adaptation.
The two spatial branches share VSSBlock weights (Siamese), but after warm-up the backbone is frozen and each branch receives its own independent LoRA adapter (Section~\ref{sec:lora}), enabling modality-specific adaptation despite the shared backbone.

\subsection{Phase 1: Warm-Up Training}
\label{sec:warmup}

During the first $W$ tasks (covering the initial geographic domain), the \backbonename{} backbone is fully trained to establish robust feature representations.
Training uses cross-entropy loss computed from cosine-similarity logits against per-class prototypes, with a reduced learning rate for the spectral branch ($0.1\times$ the base rate) to encourage early stabilization of spectral representations.
After the warm-up phase, the \emph{entire} backbone is frozen---including all three branches.
The spectral branch, having learned domain-invariant absorption signatures, becomes the permanent stable anchor.

The warm-up length $W$ is set to the number of tasks in the first domain (e.g., $W{=}3$ for MUUFL with 3 sub-tasks), ensuring complete initial-domain coverage before freezing.

\subsection{Asymmetric LoRA Adaptation}
\label{sec:lora}

For each task $t \geq W$, the frozen backbone is adapted through low-rank updates applied \emph{only to the two spatial branches}, leaving the spectral anchor untouched.

\paragraph{LoRA insertion.}
After each VSSBlock in the spatial branches, a low-rank output adapter modifies the block's feature map:
\begin{equation}
    h' = h + \alpha \cdot W_{\mathrm{up}} W_{\mathrm{down}} h
\label{eq:lora}
\end{equation}
where $h$ is the VSSBlock output, $W_{\mathrm{down}} \in \R^{r \times d}$ projects to a low-rank bottleneck, $W_{\mathrm{up}} \in \R^{d \times r}$ projects back (implemented as $1{\times}1$ convolutions for spatial feature maps), $r$ is the rank, and $\alpha{=}1.0$ is a fixed scaling factor.
Although the two spatial branches share frozen VSSBlock weights, each branch has its own independent set of LoRA adapters, enabling modality-specific adaptation.

\paragraph{Asymmetric rank allocation.}
The two spatial branches receive different LoRA ranks:
\begin{equation}
    r_{\mathrm{HSI}} = r_0, \quad r_{\mathrm{LiDAR}} = 2 r_0
\label{eq:rank}
\end{equation}
The LiDAR-spatial branch receives $2\times$ the rank because it requires a higher-dimensional adaptation subspace: as the sole carrier of elevation information with no spectral branch backup, LiDAR adaptation spans more diverse directions than HSI-spatial, whose information is partially redundant with the frozen spectral anchor.
This is confirmed by the rank allocation comparison (Table~\ref{tab:rank_allocation}), which shows that our asymmetric allocation outperforms symmetric allocation at the same total rank budget.

\paragraph{Domain-conditioned reuse.}
Rather than allocating a separate adapter per task, we share adapters across tasks within the same geographic domain.
Specifically, tasks from the same domain share a single adapter set.
For the 9-task benchmark (3~domains, 3~tasks each), this yields 2 adapter sets for 6 post-warmup tasks, reducing parameter count by $3\times$ over per-task allocation while acting as an implicit domain-level regularizer.

\paragraph{Domain-selective inference.}
At test time, only the adapter corresponding to the query's domain is activated:
\begin{equation}
    h' = h + \alpha\, W_{\mathrm{up}}^{(d)} W_{\mathrm{down}}^{(d)}\, h
\label{eq:domain_selective}
\end{equation}
where $d$ is the domain index of the test sample and $h$ is the frozen backbone feature.
Each domain's adapter is calibrated for its own distributional characteristics; activating only the relevant adapter avoids interference from adapters trained on unrelated domains.
During \emph{training}, when learning a new domain $k$, a fresh adapter is allocated and fine-tuned on the new domain's data; when learning additional tasks within the same domain, the existing domain adapter is unfrozen and continues training.
During \emph{evaluation}, the domain identity is known because test sets are physically separate per geographic region (see SHINE discussion in Section~\ref{sec:shine}).

\subsection{SHINE: Per-Domain Whitening Normalization}
\label{sec:shine}

Cross-domain prototype matching suffers from distributional misalignment: prototypes computed from one domain's training features may not be directly comparable to query features from another domain due to differences in feature statistics (mean, variance).
SHINE addresses this by applying per-domain, per-branch whitening normalization at both prototype construction and inference time.

\paragraph{Construction.}
After training on each task, for each branch $b$ and each domain $d$, we compute the running mean $\mu_d^b$ and standard deviation $\sigma_d^b$ from the training features of that domain.
Class prototypes are stored \emph{after} Z-score normalization:
\begin{equation}
    \bar{f}_b = \frac{f_b - \mu_d^b}{\sigma_d^b + \epsilon}
\label{eq:shine_norm}
\end{equation}
where $d$ is the domain from which the features originate.
When the sample count for a domain--branch pair is below a threshold ($n < 100$), the normalized features are blended with the original features using $\alpha = n/100$ to avoid unreliable statistics from small samples.

\paragraph{Inference.}
In our benchmark experiments, we use a \emph{domain-aware} mode: the query's domain identity is known because test sets are physically separate per domain, and the query is normalized with its domain's statistics.
In open deployment where domain labels are unavailable, a \emph{domain-agnostic} extension could normalize the query using each candidate domain's statistics in turn and select the highest-scoring class across all domain--prototype pairs.
We leave validation of this domain-agnostic mode to future work; the current evaluation uses domain-aware normalization throughout.

\paragraph{Deployment protocol.}
At training time, per-domain statistics $(\mu_d^b, \sigma_d^b)$ are computed from training features and stored alongside class prototypes (4.5\,KB for 3 domains $\times$ 3 branches $\times$ mean/std $\times$ $d{=}64$).
This post-hoc normalization eliminates the first-order distributional gap (mean and variance shift) between domains without introducing any learnable parameters.
SHINE is an integral component of the \methodname{} framework (analogous to herding selection in iCaRL) and is not applied to baselines.

\subsection{DCR: Drift-Compensated Reconstruction}
\label{sec:dcr}

When spatial adapters are trained on a new domain, the spatial features of old-class prototypes become stale: they were computed with the previous adapter state, which no longer matches the current accumulated adapter configuration.
DCR ensures that old-class prototypes remain aligned with the current accumulated adapter configuration.

\paragraph{Principle.}
After training the adapter for a new domain, old-class prototypes from the \emph{same} domain must reflect the updated adapter state.
For each spatial-branch prototype $\mu_c^b$ ($b \in \{\mathrm{hsi}, \mathrm{lid}\}$) of a class $c$ belonging to domain $d$, the corrected prototype under domain-selective inference is:
\begin{equation}
    \hat{\mu}_c^b = (I + \alpha\, W_{\mathrm{up}}^{(d)} W_{\mathrm{down}}^{(d)})\, \mu_c^b
\label{eq:dcr}
\end{equation}
This ensures prototypes and test features are aligned under the same domain adapter.
Spectral-branch prototypes require no correction since the spectral branch is frozen.
Lemma~1 proves this correction is exact when LoRA acts as a linear additive perturbation at the feature level; in practice, LoRA is inserted within multi-block branches where inter-block nonlinearities introduce a small approximation gap.

\paragraph{Implementation.}
In our experiments, we implement DCR implicitly: after each task, features for all seen classes are re-extracted through the current model (with the domain-specific adapter for each class's domain), and prototypes are recomputed from these updated features.
For a strict deployment scenario without old-data access, Eq.~\eqref{eq:dcr} provides an approximate correction that can be applied directly to stored prototypes; the approximation quality depends on the magnitude of the LoRA update relative to the branch's nonlinearities.

\paragraph{Ablation.}
The ``w/o DCR pipeline'' ablation (Table~\ref{tab:ablation}) disables the full domain-conditioned pipeline: adapter reuse, domain-aware distillation, and prototype correction.
Because these components are architecturally coupled, we report the combined effect ($-1.6$\,pp); see Section~\ref{sec:ablation} for discussion.

\subsection{Multi-Centroid Prototype Classification}
\label{sec:classifier}

Each class is represented by $K{=}2$ centroids per branch, obtained by $k$-means clustering on the training features after each task.
For a test sample $x$ with SHINE-normalized branch features $\{\bar{f}_b\}$, the per-branch score aggregates cosine similarities to the $K$ centroids via log-sum-exp:
\begin{equation}
    s_b(c \mid x) = \log \sum_{k=1}^{K} \pi_{c,k}^b \cdot \exp\!\left(\tau \cdot \cosim(\bar{f}_b,\; \bar{\mu}_{c,k}^b)\right)
\label{eq:multi_centroid_score}
\end{equation}
where $\pi_{c,k}^b = n_{c,k}^b / n_c^b$ is the proportion of class-$c$ training samples assigned to centroid $k$ in branch $b$, $\bar{\mu}_{c,k}^b$ is the SHINE-normalized centroid, and $\tau$ is a temperature parameter.
The final class score aggregates across branches:
\begin{equation}
    s(c \mid x) = \sum_{b \in \{\mathrm{spec}, \mathrm{hsi}, \mathrm{lid}\}} s_b(c \mid x)
\label{eq:ncm}
\end{equation}
The predicted class is $\hat{y} = \argmax_{c \in \mathcal{C}_{1:t}} s(c \mid x)$.

The multi-centroid representation captures intra-class heterogeneity (e.g., a land-cover class spanning two spectral sub-populations), improving discrimination over a single-centroid nearest-class-mean.
This prototype-based classifier avoids the bias toward recent classes that plagues fully-connected heads in CIL settings and requires no additional learnable parameters beyond the stored centroids and mixing weights.

\subsection{GMM-Weighted Orthogonal Projection}
\label{sec:gmm_ortho}

A central challenge in exemplar-free CIL is training LoRA adapters on new-class data \emph{without} degrading old-class features.
When the adapter for domain $d$ is trained on task $t$'s new classes, the low-rank update $\Delta W_b = \alpha\, W_{\mathrm{up}}^{(b)} W_{\mathrm{down}}^{(b)}$ modifies the spatial feature space.
Without constraint, this update can displace old-class prototypes, causing catastrophic forgetting (a na\"ive exemplar-free variant degrades to $33.9\%$; see Section~\ref{sec:limitations}).

We propose a \emph{GMM-weighted orthogonal projection} loss that constrains LoRA updates to preserve old-class prototype directions, using only stored GMM statistics---no old-class data access required.

\paragraph{Formulation.}
For each old class $c$ in the current domain with GMM centroids $\{\mu_{c,k}^b\}_{k=1}^{K}$ and per-dimension standard deviations $\{\sigma_{c,k}^b\}_{k=1}^{K}$ in spatial branch $b \in \{\mathrm{hsi}, \mathrm{lid}\}$, the LoRA perturbation applied to centroid $\mu_{c,k}^b$ is:
\begin{equation}
    \delta_{c,k}^b = \alpha\, W_{\mathrm{up}}^{(b)} W_{\mathrm{down}}^{(b)}\, \mu_{c,k}^b
\label{eq:gmm_perturbation}
\end{equation}
This perturbation should be small, especially along dimensions where the class distribution is tight (low variance = high discriminability).
The importance weight for dimension $j$ is inversely proportional to the variance:
\begin{equation}
    w_{c,k,j}^b = \frac{1}{(\sigma_{c,k,j}^b)^2 + \epsilon}
\label{eq:importance_weight}
\end{equation}
The GMM orthogonal projection loss penalizes the importance-weighted perturbation across all old-class centroids:
\begin{equation}
    \mathcal{L}_{\mathrm{GMM}} = \frac{1}{|\mathcal{S}|} \sum_{c \in \mathcal{C}_{\mathrm{old}}^d} \sum_{k=1}^{K} \sum_{b} \sum_{j=1}^{d} w_{c,k,j}^b \cdot (\delta_{c,k,j}^b)^2
\label{eq:gmm_loss}
\end{equation}
where $\mathcal{C}_{\mathrm{old}}^d$ denotes old classes in the current domain $d$, and $|\mathcal{S}|$ is the total number of centroid--branch terms for normalization.

\paragraph{Geometric interpretation.}
Each GMM centroid defines a direction in feature space that is important for classifying that class.
The variance-weighted penalty ensures that LoRA updates are suppressed along \emph{discriminative} dimensions (where $\sigma^2$ is small and the class is tightly clustered) while allowing adaptation along \emph{non-discriminative} dimensions (where $\sigma^2$ is large and variation is tolerable).
With $K{=}2$ centroids per class and $|\mathcal{C}_{\mathrm{old}}^d|$ old classes in the current domain, the protected subspace spans at most $2 \times |\mathcal{C}_{\mathrm{old}}^d|$ directions per branch in the $d$-dimensional feature space, leaving ample room for new-class adaptation.

\paragraph{Why same-domain only.}
Cross-domain old classes are already protected by a separate mechanism: domain-selective inference (Section~\ref{sec:lora}) ensures that each domain's test features pass through only that domain's adapter.
Combined with frozen backbone BatchNorm (Section~\ref{sec:bn_freeze}), cross-domain features are \emph{exactly invariant} to new-domain LoRA training.
Therefore, $\mathcal{L}_{\mathrm{GMM}}$ only needs to constrain same-domain old-class directions, keeping the protected subspace compact.

\paragraph{Connection to drift asymmetry.}
The GMM orthogonal projection directly exploits the drift asymmetry:
the spectral branch is frozen and requires no protection (its prototypes never change), so $\mathcal{L}_{\mathrm{GMM}}$ operates only on the spatial branches ($b \in \{\mathrm{hsi}, \mathrm{lid}\}$).
The asymmetric rank allocation (Section~\ref{sec:lora}) means the LiDAR branch has more adaptation capacity ($r_{\mathrm{lid}} = 2r_0$), and correspondingly more protected directions from old-class centroids.

\subsection{Exemplar-Free Training Protocol}
\label{sec:ef_protocol}

Combining the above components, \methodname{} achieves true exemplar-free training: each task uses \emph{only its own new-class data}, with no access to any previous task's training samples.

\paragraph{Backbone BatchNorm freezing.}
\label{sec:bn_freeze}
During LoRA training, the backbone is frozen (\texttt{requires\_grad=False}), but a subtle issue arises: calling \texttt{model.train()} switches BatchNorm layers to training mode, causing their running statistics to be updated by the current domain's data.
This pollutes old-domain features at inference time, as BatchNorm running means and variances no longer reflect the old domain's distribution.
We address this by keeping the backbone in evaluation mode during LoRA training:
\begin{equation}
    \texttt{model.train();\quad model.backbone.eval()}
\label{eq:bn_freeze}
\end{equation}
This ensures BatchNorm running statistics remain frozen, making cross-domain feature extraction \emph{deterministic} regardless of which domain was last trained.
We verified this empirically: after the fix, non-current-domain feature statistics are \emph{exactly unchanged} across all tasks (confirmed via A/B comparison).

\paragraph{New-class-only LoRA training.}
For each post-warmup task $t$ in domain $d$, the LoRA adapter is trained using only the new classes introduced by task $t$:
\begin{equation}
    \mathcal{D}_{\mathrm{train}}^{(t)} = \{(x_i, y_i) : y_i \in \mathcal{C}_t\}
\end{equation}
where $\mathcal{C}_t$ is the set of new classes in task $t$.
Old-class preservation is handled entirely through the GMM orthogonal projection (Eq.~\ref{eq:gmm_loss}) and domain-aware feature distillation ($\mathcal{L}_{\mathrm{DKD}}$), neither of which requires old-class data.

\paragraph{Current-domain-only prototype recomputation.}
After training, prototypes and SHINE statistics are recomputed only for the \emph{current domain}.
Cross-domain prototypes remain cached from their last computation and are still valid because:
(i)~the backbone is frozen,
(ii)~BatchNorm running statistics are frozen (Eq.~\ref{eq:bn_freeze}), and
(iii)~domain-selective inference ensures cross-domain features pass through unchanged adapters.
This eliminates all cross-domain data access during the training pipeline.

\paragraph{Data access summary.}
Table~\ref{tab:data_access} summarizes the data access pattern.
At every stage, \methodname{} accesses only the \emph{current task's new-class training data}---the same data that every CIL method must see.
No replay buffer is maintained, and no old-task data is accessed at any point.

\begin{table}[t]
\centering
\caption{Data access per task. \methodname{} only accesses current-task new-class data at every stage.}
\label{tab:data_access}
\small
\begin{tabular}{lll}
\toprule
\textbf{Stage} & \textbf{Data accessed} & \textbf{Old data?} \\
\midrule
LoRA training & Current task new classes & No \\
GMM ortho loss & Stored GMM statistics & No \\
DKD loss & Teacher model snapshot & No \\
Prototype update & Current task new classes & No \\
SHINE update & Current task new classes & No \\
Cross-domain & Cached (unchanged) & No \\
\bottomrule
\end{tabular}
\end{table}

\subsection{Overall Training and Inference}
\label{sec:overall}

\paragraph{Training objective.}
The total loss for task $t \geq W$ combines cross-entropy on new-class data with two regularization terms that protect old classes without requiring old-class data:
\begin{equation}
    \mathcal{L} = \Lce + \lambda_{\mathrm{GMM}} \cdot \mathcal{L}_{\mathrm{GMM}} + \lambda_{\mathrm{DKD}} \cdot \mathcal{L}_{\mathrm{DKD}}
\label{eq:total_loss}
\end{equation}
where $\Lce$ is cross-entropy on the current task's new classes, $\mathcal{L}_{\mathrm{GMM}}$ is the GMM-weighted orthogonal projection loss (Eq.~\ref{eq:gmm_loss}) that constrains LoRA updates to preserve old-class prototype directions, and $\mathcal{L}_{\mathrm{DKD}}$ is a domain-aware feature distillation loss from a teacher snapshot of the previous domain.
DCR prototype correction (Eq.~\ref{eq:dcr}) is applied \emph{after} training for classes whose data was accessed during the current task.

\paragraph{Inference pipeline.}
Given a test sample $x$ from domain $d$:
\begin{enumerate}
    \item Extract $\fspec$, $\fhsi$, $\flid$ from the frozen backbone (BatchNorm in eval mode).
    \item Apply domain-selective LoRA: activate only the adapters associated with domain $d$.
    \item Apply SHINE normalization (Eq.~\ref{eq:shine_norm}) using domain $d$'s statistics.
    \item Compute multi-centroid scores (Eq.~\ref{eq:ncm}) against all stored prototypes.
    \item Predict $\hat{y} = \argmax_c\, s(c \mid x)$.
\end{enumerate}

\paragraph{Parameter budget.}
Each domain adapter pair adds ${\sim}6$K trainable parameters (${\sim}0.20\%$ of the frozen \backbonename{} backbone), with multi-centroid prototypes (2 centroids $\times$ 32 classes $\times$ 3 branches $\times$ 64 dims $\approx$ 48\,KB) and SHINE statistics requiring negligible additional storage---orders of magnitude below the cost of storing exemplar images.
